\subsection{Descent of a primitive-zero factor through an abelian total cover}

\begin{claim}{A.C2.23}{Proved}
Let $k$ be algebraically closed of characteristic $p>0$,
let $C/k$ be a smooth proper connected curve, and let
$W\to C$ be a connected finite \'{e}tale Galois cover
with finite abelian group $A$, of odd order prime to $p$.
Let $H\subset A$ be a nontrivial cyclic subgroup of
order $n$, let $h$ generate $H$, and put $V=W/H$.
Suppose that
\[
 P_H=\operatorname{im}\left(
       \prod_{\ell\mid n}(1-h^{n/\ell}):J_W\longrightarrow J_W
                         \right)
\]
has $p$-rank zero. Then an intermediate quotient
$D=W/K$ is a cyclic \'{e}tale cover of $C$, of odd
order $N>1$ prime to $p$, whose primitive conductor-$N$
factor has $p$-rank zero. One can choose $N$ divisible
by $n$.

Thus the existence of a cyclic primitive-zero factor
over an intermediate base inside an abelian total cover
is equivalent to the existence of a cyclic
primitive-zero cover of the original curve.
For genus at least two over a finitely generated field,
this condition implies A1 by A.C2.20.
The assertion concerns an abelian \emph{total} cover
$W/C$; a sequence of abelian covers can have a
nonabelian normal closure over $C$ and is not
covered by this equivalence.
\end{claim}
\begin{proof}
\proofparagraph{Primitive characters and geometric Tate modules.}
Write
\[
 \mathcal T_W=T_p(J_W)\otimes_{\mathbf Z_p}
                              \overline{\mathbf Q}_p.
\]
The geometric $p$-adic Tate module has rank equal
to the $p$-rank. Since $A$ has order prime to $p$,
its action is semisimple; since it is abelian,
\[
 \mathcal T_W=\bigoplus_{\psi\in A^\vee}\mathcal T_{W,\psi},
 \qquad
 A^\vee=\operatorname{Hom}(A,\overline{\mathbf Q}_p^\times).
\]
On the $\psi$-space the operator defining $P_H$
has scalar
\[
              \prod_{\ell\mid n}
                   \bigl(1-\psi(h)^{n/\ell}\bigr).
\]
This scalar is nonzero precisely when $\psi|_H$
is faithful. Indeed, if its order $m$ properly
divides $n$, some prime $\ell$ satisfies
$v_\ell(m)<v_\ell(n)$, whence $m\mid n/\ell$.
If its order is $n$, no factor vanishes.

The image on a rational geometric Tate module of
a homomorphism of abelian varieties is the Tate
module of its image abelian subvariety. This
follows by factoring through the connected
kernel and the resulting isogeny; isogenies
become isomorphisms on rational $p$-adic Tate
modules. Consequently
\[
 \mathcal T_{W,\psi}=0
          \quad\text{if }\psi|_H\text{ is faithful}.
                                                        \tag{AT.1}
\]
Here zero $p$-rank is used, not merely a
nonzero Frobenius kernel.

\proofparagraph{Extension of a faithful subgroup character.}
Choose a faithful character
$\theta:H\to\overline{\mathbf Q}_p^\times$.
It extends to a character $\chi$ of $A$.
Here is the elementary extension argument.
Suppose a character is defined on a subgroup $B$,
choose $a\in A$, and let $r$ be the order of
$a+B$ in $A/B$. Choose $z$ with
$z^r=\theta(ra)$. The formula
\[
               b+ja\longmapsto\theta(b)z^j
\]
defines a character of $B+\langle a\rangle$.
An equality of two such expressions has
$j-j'$ divisible by $r$, and the root relation
gives precisely the necessary equality of
values. If $a$ has order $s$, then $r\mid s$
and $z^s=\theta(sa)=1$, so the image remains
finite. Finitely many enlargements give $\chi$.

Let $N$ be the order of $\chi$, put
$K=\ker\chi$, and set $D=W/K$.
Then $A/K$ is cyclic of order $N$ and
$D\to C$ is a connected cyclic \'{e}tale cover.
Since $\chi|_H$ has order $n$, one has $n\mid N$.
For every $u\in(\mathbf Z/N)^\times$,
\[
                    (\chi^u)|_H=\theta^u
\]
is faithful: $n\mid N$ implies $\gcd(u,n)=1$.
Equation (AT.1) therefore gives
\[
       \mathcal T_{W,\chi^u}=0
                 \quad(u\in(\mathbf Z/N)^\times).
                                                        \tag{AT.2}
\]

\proofparagraph{Norm and pullback for the cyclic quotient.}
For $q:W\to D$, pullback and norm satisfy
\[
 q_*q^*=[|K|],\qquad
 q^*q_*=\sum_{\gamma\in K}\gamma^*.
\]
Since $|K|$ is prime to $p$, they identify
$T_p(J_D)$ with $T_p(J_W)^K$, already
integrally. For an invariant vector $v$, the
inverse comparison is explicit:
\[
                   v=q^*\bigl(|K|^{-1}q_*v\bigr).
\]
The $A$-characters trivial on $K$ are exactly
the powers of $\chi$. The faithful characters
of $A/K$ are exactly the primitive powers
$\chi^u$, $u\in(\mathbf Z/N)^\times$.
Their Tate spaces vanish by (AT.2), so the
primitive conductor-$N$ factor of $J_D$ has
$p$-rank zero.
Dualizing conventions may invert character
indices; faithfulness and primitivity are
unchanged.

This proves descent. The converse existence
implication takes $W=D$, $A=H$, and $V=C$.
All finite geometric covering data descend
after a finite extension of a field of
definition. Thus A.C2.20 supplies the stated
finitely generated-field consequence, including
its forms and open-subcurve variants.
\end{proof}

For genus two in characteristic two, this
explains an exact obstruction to an abelian
base-change strategy. Passing to an
intermediate base groups the original
character spaces according to their
restrictions to the new deck subgroup.
One theta hit need not kill all original
characters extending a faithful character
upstairs. If every such contribution does
vanish, the proof above already produces a
cyclic quotient over the original curve
with all primitive doubling orbits killed.
In particular A.C2.14's finite-prime-set
avoidance cannot be circumvented by taking
intermediate cyclic stages of an abelian
total cover supported on the same prime set.

This is a limitation of that construction,
not a counterexample to A1 or a
nonexistence assertion for nonabelian
normal closures. The separate source audit
\path{audit/A_nonabelian_Raynaud_gap.md}
records why the opened higher-rank
nonordinary-cover theorem does not itself
supply zero stable multiplicities or a
sparse integral annihilator.
The finite checks in
\path{code/BC_abelian_descent_and_support_checks.py}
verify 31,006 character assertions and
an explicitly nongeometric Heisenberg
module illustrating the distinction
between representation support and
multiplicity.
